\subsection{Reduced Averaging Impact}
\label{sec:discussion-averaging}

The 16-fold reduction in the averaging protocol (16 vs.\ 256 averages) carries substantial practical implications for experimental workflows. In terms of acquisition speed, the factor-of-16 acceleration directly translates to a $16\times$ throughput increase, enabling high-throughput screening and in-situ inspection regimes that would otherwise be prohibitively time-consuming. Complementing this is a $16\times$ reduction in cumulative laser fluence on the specimen surface, which is critical when working with sensitive coatings, fragile biological tissues, or photo-sensitive materials where thermal accumulation or photobleaching degrades sample integrity. The method therefore addresses not only efficiency but also sample preservation constraints inherent to optical imaging pipelines.

Quantitatively, the efficiency gain is captured by the equivalence PMDF-Net$+$16 averages $\approx$ averaging$+$256 averages, as demonstrated in Section~\ref{sec:results}. This near-parity means that the network effectively learns a denoising function that replaces what would otherwise require a 16-fold increase in photon budget. The practical consequence is that researchers can operate within tighter acquisition budgets without sacrificing image quality, or equivalently, push toward higher quality at a given dose.

The trade-off inherent to this approach is the requirement for paired low/high-averaging training data. The network must learn the residual mapping between under-averaged and fully-averaged acquisitions rather than inferring it from statistics alone. Fortunatey, the present study includes this paired training corpus (Section~\ref{sec:method-inference}), ensuring the reduced-averaging protocol is fully validated. Institutions adopting this framework must therefore establish equivalent paired datasets for their specific specimens and imaging conditions. The reduced averaging protocol thus shifts the burden from photon collection at inference time to data curation at training time, a trade-off that favors practical experimental workflows where acquisition time and dose are limiting factors.